\section{The \texttt{Mirror\_parabolic} McXtrace Component}
Idealized parabolic mirror (in XZ)

\subsection*{Identification}
\begin{itemize}
  \item \textbf{Author:} Erik Knudsen
  \item \textbf{Origin:} Risoe
  \item \textbf{Date:} Feb 11, 2010
\end{itemize}

\subsection*{Description}
Takes a reflectivity as input and reflects rays in a ideal geometry parabolic mirror. The mirror is positioned in the zx-plane curving towards positive y. I.e. the focal point is (0,0,f(a,b)) The geometry of the paraboloid is governed by the equation: y = x\textasciicircum{}2 / a\textasciicircum{}2 + z\textasciicircum{}2 / b\textasciicircum{}2 Hence, the focal length for the 'x' curve is f=a\textasciicircum{}2 / 4, and analogous for z.

Example: Mirror\_parabolic(R0=1, a=1, b=0, xwidth=0.02, yheight=0, zdepth=0.05)

\subsection*{Input parameters}
Parameters in \textbf{boldface} are required; the others are optional.

\begin{longtable}{p{0.22\textwidth}p{0.12\textwidth}p{0.46\textwidth}p{0.14\textwidth}}
\toprule
\textbf{Name} & \textbf{Unit} & \textbf{Description} & \textbf{Default} \\
\midrule
\endhead
R0 & 1 & Reflectivity of mirror. & 1 \\
a & sqrt(m) & Transverse curvature scale, if zero - the mirror is flat along x. & 1 \\
b & sqrt(m) & Longitudinal curvature scale, if zero, flat along z. & 1 \\
xwidth & m & Width of mirror. & 0.1 \\
zdepth & m & Length of mirror. & 0.1 \\
yheight & m & Thickness of mirror. If 0 (the default) the mirror is mathemticlly thin. Only has an effect for hitting the mirror from the side. & 0 \\
focusx & m & Transverse focal length along X. Sets a. & 0 \\
focusz & m & Longitudinal focal length along Z. Sets b. & 0 \\
radius & m & Focal length. Sets focusx and focusz. & 0 \\
\bottomrule
\end{longtable}

\subsection*{Links}
\begin{itemize}
  \item Component source code found in file \texttt{Mirror\_parabolic.comp}.
\end{itemize}
\IfFileExists{optics/Mirror_parabolic_static.tex}{\input{optics/Mirror_parabolic_static.tex}}{}